\documentclass[french]{article}

\usepackage[utf8]{inputenc}
\usepackage[T1]{fontenc}
\usepackage{lmodern}
\usepackage{geometry}
\usepackage{graphicx}
\usepackage{babel}
\usepackage{amsmath}
\usepackage{amsthm}
\usepackage{amssymb}
\usepackage{hyperref}
\usepackage{stmaryrd}
\usepackage{enumitem}
\usepackage[most]{tcolorbox}
\usepackage{dsfont}
\usepackage{multirow}
\usepackage{etoc}
\usepackage{titlesec}
\usepackage{esint}
\usepackage{cancel}
\usepackage{mathdots}
\usepackage{indentfirst}

\setlength{\parindent}{0pt}

\setlist[itemize]{label=$\bullet$}
\setlist{before=\leavevmode, leftmargin=*}

\renewcommand{\P}{\mathbb{P}}
\newcommand{\N}{\mathbb{N}}
\newcommand{\Z}{\mathbb{Z}}
\newcommand{\Q}{\mathbb{Q}}
\newcommand{\R}{\mathbb{R}}
\newcommand{\C}{\mathbb{C}}
\newcommand{\K}{\mathbb{K}}
\renewcommand{\L}{\mathbb{L}}
\newcommand{\U}{\mathbb{U}}
\newcommand{\st}{^*}
\newcommand{\tq}{\middle|}
\newcommand{\inv}{^{-1}}
\newcommand{\E}{\mathbb{E}}
\newcommand{\F}{\mathbb{F}}
\newcommand{\V}{\mathbb{V}}
\renewcommand{\ker}{\text{Ker}}
\newcommand{\im}{\text{Im}}
\newcommand{\card}{\text{Card}}
\newcommand{\Tr}{\text{Tr}}
\newcommand{\Vect}{\text{Vect}}
\newcommand{\rg}{\text{rg}}
\newcommand{\vertiii}[1]{{\left\vert\kern-0.25ex\left\vert\kern-0.25ex\left\vert #1 \right\vert\kern-0.25ex\right\vert\kern-0.25ex\right\vert}}
\renewcommand{\o}{\text{o}}
\renewcommand{\O}{\text{O}}
\renewcommand{\d}{\text{d}}
\newcommand{\Id}{\text{Id}}


\theoremstyle{plain}
\newtheorem*{ind}{Indication}

\theoremstyle{definition}

\newtheorem*{ex}{Exemple}
\newtheorem*{rmq}{Remarque}
\newtheorem*{notation}{Notation}
\newtheorem*{rappel}{Rappel}
\newtheorem*{terminologie}{Terminologie}
\newtheorem*{attention}{Attention}
\newtheorem*{conséquence}{Conséquence}

\theoremstyle{remark}
\newtheorem*{app}{Application}

\newtcolorbox[auto counter]{dfn}[1][]{
	enhanced,
	colback=white,
	colframe=black,
	sharp corners,
	boxrule=0.8pt,
	fonttitle=\bfseries,
	colbacktitle=white,
	coltitle=black,
	attach boxed title to top left={yshift=-1mm,xshift=-0.5mm},
	boxed title style={size=minimal, right=2pt, sharp corners, colframe=white},
	title={Définition \thetcbcounter \ifstrempty{#1}{}{ (#1)}},
	before skip=0.5em,
	after skip=0.5em
}

\newtcolorbox[auto counter]{exo}[1][]{
	enhanced,
	arc=1mm,
	colback=white,
	colframe=black,
	coltitle=black,
	fonttitle=\bfseries,
	boxrule=0.8pt,
	shadow={1.2pt}{-1.2pt}{0pt}{black!50},
	attach title to upper={\ },
	title={Exercice \thetcbcounter\ifblank{#1}{}{ #1}},
	before skip=0.5em,
	after skip=0.5em,
	left=2mm, right=2mm, top=1mm, bottom=1mm
}

\newtcolorbox{met}[1][]{
	enhanced,
	arc=1mm,
	colback=white,
	colframe=black,
	coltitle=black,
	fonttitle=\bfseries,
	boxrule=0.8pt,
	shadow={1.2pt}{-1.2pt}{0pt}{black!50},
	attach title to upper={\ },
	title={Point méthode \ifblank{#1}{}{#1}},
	before skip=1em,
	after skip=1em,
	left=2mm, right=2mm, top=1mm, bottom=1mm
}

\newcounter{thmcount}

\newtcolorbox[use counter=thmcount]{thm}[1][]{
	enhanced,
	colback=white,
	colframe=black,
	sharp corners,
	left skip=0mm,
	boxrule=1.3pt,
	fonttitle=\bfseries,
	colbacktitle=white,
	coltitle=black,
	attach boxed title to top left={yshift=-1mm,xshift=-0.5mm},
	boxed title style={size=minimal, right=2pt, sharp corners, colframe=white},
	title={Théorème \thethmcount \ifstrempty{#1}{}{ (#1)}},
	before skip=0.5em,
	after skip=0.5em
}

\newtcolorbox[use counter=thmcount]{prop}[1][]{
	enhanced,
	colback=white,
	colframe=black,
	sharp corners,
	boxrule=0.8pt,
	fonttitle=\bfseries,
	colbacktitle=white,
	coltitle=black,
	attach boxed title to top left={yshift=-1mm,xshift=-0.5mm},
	boxed title style={size=minimal, right=2pt, sharp corners, colframe=white},
	title={Proposition \thethmcount \ifstrempty{#1}{}{ (#1)}},
	before skip=0.5em,
	after skip=0.5em
}

\newtcolorbox[use counter=thmcount]{cor}[1][]{
	enhanced,
	colback=white,
	colframe=black,
	sharp corners,
	boxrule=0.8pt,
	fonttitle=\bfseries,
	colbacktitle=white,
	coltitle=black,
	attach boxed title to top left={yshift=-1mm,xshift=-0.5mm},
	boxed title style={size=minimal, right=2pt, sharp corners, colframe=white},
	title={Corollaire \thethmcount \ifstrempty{#1}{}{ (#1)}},
	before skip=0.5em,
	after skip=0.5em
}

\newtcolorbox[use counter=thmcount]{lem}[1][]{
	enhanced,
	colback=white,
	colframe=black,
	sharp corners,
	boxrule=0.5pt,
	fonttitle=\bfseries,
	colbacktitle=white,
	coltitle=black,
	attach boxed title to top left={yshift=-1mm,xshift=-0.5mm},
	boxed title style={size=minimal, right=2pt, sharp corners, colframe=white},
	title={Lemme \thethmcount \ifstrempty{#1}{}{ (#1)}},
	before skip=0.5em,
	after skip=0.5em
}

\title{Compacité, connexité, dimension finie}
\date{}

\begin{document}
\tableofcontents

\newpage
\maketitle

Dans tout ce chapitre, $E$ désigne un espace métrique, par exemple un $\K$-espace vectoriel normé où $\K$ désigne $\R$ ou $\C$.\\
L'un des objectifs de ce chapitre est de généraliser les deux résultats principaux concernant la continuité des fonctions d'une variable réelle :
\begin{itemize}
	\item le théorème des valeurs intermédiaires ;
	\item le théorème des bornes atteintes ;
	\item le théorème de Heine.
\end{itemize}
Nous finirons par voir (enfin !) les conséquences sur les espaces vectoriels de dimension finie des nouvelles notions introduites.

\section{Compacité}%intéressant pour les arguments séquentiels

\subsection{Parties compactes}%généralisation intéressante de la notion de partie finie

\begin{dfn}
	On dit qu'une partie $A\subset E$ est \textbf{compacte} si toute suite de $A$ admet une valeur 'adhérence dans $A$, c'est-à-dire si de toute suite de $A$ on peut extraire une suite convergeant dans $A$.
\end{dfn}

\begin{terminologie}
	On dit parfois qu'il s'agit de la compacité au sens de Bolzano-Weierstrass.
\end{terminologie}

\begin{rmq}
	\begin{enumerate}
		\item La compacité est une notion topologique puisqu'elle ne fait intervenir que la convergence de suites. \\
		En particulier, si $E$ est un espace vectoriel normé, une partie $A$ de $E$ compacte pour une norme $N$ est compacte pour tout norme équivalente à $N$.
		\item La compacité d'une partie $A$ est \textit{intrinsèque} à $A$ : elle ne dépend que de la topologie de $A$ induite par la distance de $E$.
	\end{enumerate}
\end{rmq}

\begin{ex}
	L'ensemble vide, les singletons et les parties finies sont des compacts.
\end{ex}

\begin{prop}
	Toute partie compacte de $E$ est fermée dans $E$.
\end{prop}

\begin{prop}
	Toute partie compacte de $E$ est bornée.
\end{prop}

\begin{attention}
	Les parties compactes d'un espace métrique sont donc fermées et bornées. Mais la réciproque est fausse en général, sauf pour les espaces vectoriels de dimension finie.
\end{attention}

\begin{ex}
	Dans $\R[X]$ muni de la norme $\lVert\cdot\rVert_\infty$, la suite $X^n$ est bornée et vérifie $\lVert X^p-X^q\rVert_\infty=1$ pour $p\ne q$, donc n'admet pas de sous-suite convergente. La sphère unité est donc un fermé borné non compact.
\end{ex}

\begin{met}
	Comme on l'a fait dans l'exemple précédent, une méthode classique pour montrer qu'une partie (non fermée bornée) est non compacte est d'essayer d'en exhiber une suite $(u_n)$ vérifiant $\forall p\ne q,\ d(u_p,u_q)\geqslant\alpha$, où $\alpha$ est un réel strictement positif.
\end{met}

\begin{prop}%prop
	Toute partie fermée dans un compact est compacte.
\end{prop}

\begin{ex}%un exemple de partie compacte
	Soit $(x_p)$ une suite de $E$ convergeant vers $l\in E$. Alors $A=\{x_p\ \mid\ p\in\N\}\cup\{l\}$ est un compact.
\end{ex}

\begin{prop}%prop
	Un produit de deux compacts (et donc par récurrence d'un nombre fini de compacts) est compact pour la distance produit.
\end{prop}
"Ah bah en fait c'était trivial !" (barre sa preuve)

\subsection{Parties compactes de $\K$ et de $\K^n$}

\begin{thm}[Théorème de Bolzano-Weierstrass]%thm
	Les parties compactes de $\K^n$ sont exactement les parties fermées bornées.
\end{thm}

\begin{rmq}
	Les normes classiques $\lVert\cdot\rVert_1$ et $\lVert\cdot\rVert_2$ étant équivalentes à la norme $\lVert\cdot\rVert_\infty$, le résultat précédent vaut pour chacune de ces normes, et même pour toutes les normes puisqu'elles sont équivalentes.
\end{rmq}

\subsection{Suites convergentes dans un compact}

"Je n'ai fait que copier l'énoncé idiot du programme."

\begin{thm}%thm
	Une suite à valeurs dans un compact si, et seulement si, elle n'a qu'une seul valeur d'adhérence.
\end{thm}

\begin{met}
	Pour montrer qu'une suite $(u_n)$ à valeurs dans un compact converge vers $l$, on peut montrer que toute valeur d'adhérence de $(u_n)$ est égale à $l$.
\end{met}

\begin{ex}
	La réciproque d'une bijection continue d'un compact $A$ sur un espace métrique $B$ est continue sur $B$.
	%Soit $b\in B$. Soit $(y_n)\in B^\N$ qui converge vers $b$. On peut écrire $b=f(a)$ par surjectivité de $f$ et $x_n=f\inv(y_n)$. Montrons que toute valeur d'adhérence de $(x_n)$ est égale à $a$. Si une suite extraite de $(x_n)$ tend vers $l$, alors par continuité de $f$ en $l$ et unicité de la limite, on trouve $l=a$. On conclut par compacité de $A$.
\end{ex}

\begin{cor}%cor
	Une suite bornée de $\K^n$ n'ayant qu'une seule valeur d'adhérence est convergente.
\end{cor}

\begin{ex}
	Soit $\displaystyle\sum u_n$ une série numérique dont les sommes partielles sont bornées et le terme général tend vers $0$. Alors soit elle est convergente, soit ses sommes partielles admettent une infinité de valeurs d'adhérence.
\end{ex}

\subsection{Fonctions continues sur un compact}

\begin{prop}%prop
	L'image d'un compact par une application continue est un compact.
\end{prop}

\begin{ex}
	\begin{enumerate}
		\item Tout segment d'un espace vectoriel normé est compact.
		\item Toute bijection continue sur un compact admet une réciproque continue.
	\end{enumerate}
\end{ex}

\begin{attention}
	L'image réciproque d'un compact par une application continue peut très bien ne pas être compacte. Elle est tout au plus fermée. Produire un exemple.
\end{attention}

\begin{thm}[Théorème des bornes atteintes]%thm
	Toute application réelle continue définie sur un compact non vide est bornée et atteint ses bornes.
\end{thm}

\begin{ex}
	Soit $K$ une partie compacte non vide de $E$. Pour tout $a\in E$, il existe $x\in K$ tel que $d(a,x)=d(a,K)$.
\end{ex}

\begin{exo}
	Soit $K$ un compact non vide. Montrer qu'il existe deux points $x$ et $y$ de $K$ tels que $d(x,y)=\text{diam}(K)$.
\end{exo}

Exercice Moulin.EVN.14

\begin{thm}[Théorème de Heine]%thm
	Tout application continue sur un compact est uniformément continue.
\end{thm}

\begin{exo}
	Montrer que $f:[0,1]^2\to\R$ est continue si, et seulement si, l'on a les deux conditions suivantes :
	\begin{itemize}
		\item pour tout $x\in[0,1]$, l'application $\begin{array}{lrcl}
			f_x:&[0,1]&\longrightarrow&\R\\
			&y&\longmapsto&f(x,y)
		\end{array}$ est continue ;
		\item l'application $\begin{array}{lrcl}
			[0,1]&\longrightarrow&(\mathcal{C}([0,1],\R),\mathcal{N}_\infty)\\
			x&\longmapsto&f_x
		\end{array}$ est continue.
	\end{itemize}
\end{exo}

\subsection{Compacts et réunion/intersection (HP)}

Les résultats suivants ne sont pas au programme mais peuvent s'avérer utiles. \\
"Hors programme mais vraiment très utile, c'est donc pour ça que c'est hors programme."

\begin{thm}[Théorème des compacts emboîtés]
	Soit $(K_n)_{n\in\N}$ une suite décroissante de compacts tous non vides. Alors $\displaystyle\bigcap_{n\in\N}A_n\ne\emptyset$.
\end{thm}

\begin{cor}%cor
	Soit $(F_n)$ une suite de fermés d'un compact $K$. \\
	Si $\displaystyle\bigcap_{n\in\N} F_n=\emptyset$, alors il existe $N\in\N$ tel que $\displaystyle\bigcap_{n=0}^NF_n=\emptyset$.
\end{cor}

\begin{cor}%cor
	Soit $K$ une partie compacte d'un espace métrique $E$. Si $(U_n)$ est une suite d'ouverts de $E$ recouvre $K$ (c'est-à-dire si $K\subset\displaystyle\bigcup_{n\in\N}U_n$), alors $K$ est aussi recouvert par un nombre fini de $U_n$.
\end{cor}

Les deux résultats équivalents ci-dessus admettent une réciproque :

\begin{prop}%prop
	Soit $K\subset E$. Si de tout recouvrement de $K$ par une suite d'ouverts de $E$ on peut en extraire un sous-recouvrement fini, alors $K$ est compact.
\end{prop}
\begin{proof}
	Si $(x_n)$ est une suite de $K$, alors $K\cap\displaystyle\left(\bigcap_{n\in\N}\overline{\{x_p\ \mid\ p\geqslant n\}}\right)\ne\emptyset$ en utilisant la version équivalente à l'hypothèse concernant les fermés.
\end{proof}

\begin{ex}
	On peut redémontrer ainsi que si $(x_n)$ est une suite convergente de limite $l$, alors $\{x_p\ \mid\ p\in\N\}\cup\{l\}$ est un compact.
\end{ex}

Curieusement, le caractère dénombrable de la famille d'ouverts de la famille d'ouverts des deux résultats précédents est inutile :

\begin{thm}[Théorème de Borel-Lebesgue (FHP)]%thm
	Une partie $K$ d'un espace métrique est compacte si, et seulement si, elle vérifie la propriété de Borel-Lebesgue, c'est-à-dire si de tout recouvrement de $K$ par des ouverts on peut extraire un sous-recouvrement fini.
\end{thm}
\begin{proof}
	Voir dans la section suivante.
\end{proof}

\begin{ex}
	Une application localement bornée sur un compact est bornée sur ce compact. C'est par exemple le cas pour :
	\begin{itemize}
		\item les fonctions continues sur un compact ;
		\item les fonctions continues par morceaux sur un segment de $\R$ ;
		\item les fonctions réglées sur un segment, c'est-à-dire admettant une limite finie à droite et à gauche en tout point.
	\end{itemize}
\end{ex}

\subsection{Compacité dans un espace topologique (HP)}

\subsubsection*{Généralités sur la compacité dans les espaces topologiques}

Nous donnons ici des définitions plus générales de la compacité dans les espaces topologiques, et démontrons le lien avec les résultats précédents.

\begin{dfn}
	Un espace topologique $(E,\mathcal{O})$ est dit \textbf{quasi-compact} s'il vérifie la propriété de Borel-Lebesgue : de tout recouvrement de $E$ par des ouverts on peut extraire un sous-recouvrement fini.
\end{dfn}

\begin{dfn}
	Un espace topologique est dit \textbf{séparé} lorsque deux points distincts de l'espace admettent des voisinages disjoints.
\end{dfn}

\begin{ex}
	Tout espace métrique est séparé.
\end{ex}

\begin{prop}%prop
	Tout espace topologique séparé vérifie la propriété d'unicité de la limite.
\end{prop}

\begin{dfn}
	Un espace topologique est dit \textbf{compact} s'il est quasi-compact et séparé.
\end{dfn}

\begin{rmq}
	Chez les anglo-saxons, on appelle "compact" ce que nous appelons "quasi-compact". Ils précisent "compact de Haussdorff" pour parler de notre notion de compacité.
\end{rmq}

\subsubsection*{Équivalence entre les propriétés de Borel-Lebesgue et de Bolzano-Weierstrass dans le cadre métrique}

Pour démontrer l'équivalence énoncée plus haut, nous allons avoir besoin d'une définition et de deux lemmes.

\begin{dfn}
	Un espace métrique $E$ est dit \textbf{précompact} si pour tout $\varepsilon>0$, on peut recouvrir $E$ par un nombre \textit{fini} de boules de rayon $\varepsilon$.
\end{dfn}

\begin{lem}[Compact $\implies$ précompact]%lem
	Tout espace métrique compact $E$ est précompact.
\end{lem}
\begin{proof}[Preuve par la propriété de Borel-Lebesgue]
	Si $\varepsilon>0$, il suffit d'extraire un sous-recouvrement fini du recouvrement : $$E\subset\bigcup_{x\in E}B_o(x,\varepsilon).$$
\end{proof}
\begin{proof}[Preuve par la propriété de Bolzano-Weierstrass]
	Par l'absurde, on suppose qu'il existe $\varepsilon>0$ tel que $E$ ne soit jamais recouvert par un nombre fini de boules de rayon $\varepsilon$. On construit par récurrence une suite $(x_n)$ qui vérifie : $$\forall n\in\N,\ x_{n+1}\notin\bigcup_{i=0}^nB_o(x_i,\varepsilon).$$ Une telle suite ne possède pas de valeur d'adhérence, ce qui est absurde.
\end{proof}

\begin{lem}[Lemme de recouvrement]%lem
	Soit $E$ un espace métrique compact. On suppose que $(O_i)_{i\in I}$ est une famille d'ouverts recouvrant $E$. Alors : $$\exists\varepsilon >0,\ \forall x\in E,\ \exists i\in I,\ B_o(x,\varepsilon)\subset O_i.$$
\end{lem}
\begin{proof}
	On raisonne par l'absurde avec $\varepsilon=1/2^n$ et un $x_n$ correspondant. On prend $x$ une valeur d'adhérence de $(x_n)$. Il existe $i$ tel que $x\in O_i$. Mais alors, à partir d'un certain rang : $$B\left(x_{\varphi(n)},\frac{1}{2^{\varphi(n)}}\right)\subset O_i$$ ce qui est une contradiction.
\end{proof}

Nous pouvons désormais passer à la démonstration de l'équivalence entre les compacités au sens de Borel-Lebesgue et de Bolzano-Weierstrass.

\begin{thm}%thm
	Tout espace métrique compact au sens de Bolzano-Weierstrass est compact au sens de Borel-Lebesgue.
\end{thm}
\begin{proof}
	Un espace métrique est séparé, il suffit de prouver la propriété de Borel-Lebesgue. On prend le $\varepsilon>0$ obtenu avec le lemme de recouvrement, et on utilise la précompacité de $E$. On choisit finalement, grâce au lemme de recouvrement, un ouvert par boule de rayon $\varepsilon$.
\end{proof}

\begin{thm}%thm
	Tout espace métrique compact au sens de Borel-Lebesgue est compact au sens de Bolzano-Weierstrass.
\end{thm}
\begin{proof}
	On prouve d'abord un lemme. Soit $(F_n)$ une suite décroissante de fermés tous non vides. Raisonnons par l'absurde et supposons $$\bigcap_{n\in\N} F_n=\emptyset$$ Alors $$\bigcup_{n\in\N}\underbrace{E\setminus F_n}_{:=O_n}=E$$ Alors, par hypothèse, il existe une sous-famille $(O_i)_{i\in I}$ finie telle que $$\bigcup_{i\in I}O_i=E$$ Donc $$F_{\max(I)}=\bigcap_{i\in I}F_i=\emptyset$$ ce qui est absurde. Ensuite, on se donne une suite $(u_n)\in E^\N$ et on pose $$F_n=\overline{\{u_k\mid k\geqslant n\}}$$ $(F_n)$ est une suite décroissante de fermés tous non vides donc l'intersection de tous les $F_n$, qui vaut exactement l'ensemble des valeurs d'adhérence de $(u_n)$ est non vide. Donc $(u_n)$ admet une valeur d'adhérence. Donc $E$ est compact au sens de Bolzano-Weierstrass.
\end{proof}

\section{Connexité par arcs}

\subsection{Définitions}

Soit $A$ une partie de l'espace topologique $E$.

\begin{dfn}
	Soit $x$ et $y$ deux éléments de $A$. Un \textbf{chemin joignant $x$ à $y$ dans $A$} est une application continue $p$ définie sur un segment $[a,b]$ de $\R$, à valeurs dans $A$ et telle que $p(a)=x$ et $p(b)=y$. \\
	On dit alors aussi que les points $x$ et $y$ sont \textbf{reliés par un chemin dans $A$}.
\end{dfn}

\begin{ex}
	Si $E$ est un espace vectoriel normé et $A$ est une partie convexe de $E$, deux points quelconques $x$ et $y$ de $A$ sont reliés dans $A$ par le chemin $t\mapsto (1-t)x+ty$ défini sur le segment $[0,1]$.
\end{ex}

\begin{rmq}
	\begin{itemize}
		\item Lorsqu'il existe un chemin joignant $x$ à $y$ dans $A$, celui-ci peut être défini sur $[0,1]$. En effet, si $p:[a,b]\to A$ est une application continue telle que $p(a)=x$ et $p(b)=y$, alors $q:t\mapsto p((1-t)a+tb)$ est une application continue de $[0,1]$ dans $A$ telle que $q(0)=p(a)=x$ et $q(1)=p(b)=y$.
		\item De même, on peut paramétrer un chemin joignant $x$ à $y$ par n'importe quel segment non trivial de $\R$.
		\item S'il existe un chemin joignant $x$ à $y$ et un chemin joignant $y$ à $z$, alors il existe un chemin joignant $x$ à $z$ (en "collant" les chemins).
	\end{itemize}
\end{rmq}

\begin{dfn}
	On dit que $A$ est \textbf{connexe par arcs} si deux éléments quelconques de $A$ peuvent être reliés par un chemin.
\end{dfn}

\begin{prop}%prop
	Toute partie convexe d'un espace vectoriel normé est connexe par arcs.
\end{prop}

\begin{ex}
	\begin{enumerate}
		\item Tout sous-espace affine d'un espace vectoriel normé est connexe par arcs.
		\item Tout espace vectoriel normé est connexe par arcs.
		\item L'ensemble vide est connexe par arcs.
	\end{enumerate}
\end{ex}

\begin{exo}
	$\R\st$ et $\C\st$ sont-ils connexes par arcs pour la topologie usuelle ?
\end{exo}

\begin{rmq}
	\begin{itemize}
		\item Cette notion est topologique : elle ne fait intervenir que de la continuité, donc ne dépend que de la topologie de $E$.\\
		En particulier, si $E$ est un espace vectoriel  normé, on peut donc remplacer sa norme par n'importe quelle norme équivalente sans changer ses parties connexes par arcs.
		\item La connexité par arcs de $A$ est \textit{intrinsèque} à $A$ : elle ne dépend que de la topologie de $A$ induite par celle de $E$.
	\end{itemize}
\end{rmq}

\begin{dfn}
	Soit $A$ une partie d'un espace vectoriel normé $E$ et $a\in A$. On dit que $A$ est \textbf{étoilée par rapport à $a$} si : $$\forall x\in A,\ [a,x]\subset A.$$
\end{dfn}

\begin{ex}
	\begin{enumerate}
		\item Une partie d'un espace vectoriel normé est convexe si, et seulement si, elle est étoilée par rapport à chacun de ses points.
		\item L'ensemble des matrices nilpotentes, ainsi que l'ensemble des matrices non inversibles, sont étoilés par rapport à la matrice nulle.
	\end{enumerate}
\end{ex}

\begin{prop}%prop
	Toute partie d'un espace vectoriel normé étoilée par rapport à l'un de ses points est connexe par arcs.
\end{prop}

\subsection{Composantes connexes par arcs}

\begin{prop}%prop
	La relation "il existe un chemin joignant $x$ à $y$ dans $A$" est une relation d'équivalence sur $A$.\\
	Ses classes d'équivalence sont appelées \textbf{composantes connexes par arcs de $A$}.
\end{prop}

\begin{rmq}
	\begin{itemize}
		\item A nouveau, c'est une notion qui ne dépend que la topologie de $A$.
		\item $A$ est connexe par arcs si, et seulement si, elle ne possède qu'une seule composante connexe par arcs.
	\end{itemize}
\end{rmq}

\begin{ex}
	Les composantes connexes par arcs sont des connexes par arcs.
\end{ex}

\begin{exo}
	Quelles sont les composantes connexes par arcs de $\Q$ ? de $\R\st$ ? de $\C\st$ ? (pour la topologie usuelle)
\end{exo}

\begin{exo}
	Les résultats suivants sont en général faux dans un espace métrique ou un espace topologie. Voir l'exercice Moulin.Comp.39
	\begin{enumerate}
		\item Montrer que les composantes connexes par arcs d'un ouvert $U$ d'un espace vectoriel normé sont des ouverts.
		\item Sont-elles fermées dans $U$ ?
	\end{enumerate}
\end{exo}

\subsection{Opérations sur les connexes par arcs}

\paragraph{Propriétés}
\begin{itemize}
	\item Une réunion de connexes par arcs n'est pas nécessairement connexe par arcs.
	\item La réunion d'une famille de connexes par arcs d'intersection deux à deux non vide est connexe par arcs.
	\item L'intersection de deux connexes par arcs n'est pas nécessairement connexe par arcs.
	\item Un produit fini de connexes par arcs est connexe par arcs. 
\end{itemize}

\begin{ex}
	Si $X$ est une partie finie de $\C$, alors $\C\setminus X$ est connexe par arcs.
\end{ex}

\begin{exo}
	Et si $X$ est dénombrable ?
\end{exo}

\subsection{Parties connexes par arcs de $\R$}

\begin{prop}%prop
	Les parties connexes par arcs de $\R$ sont les intervalles.
\end{prop}

\begin{exo}
	Montrer que tout ouvert de $\R$ est une réunion au plus dénombrable d'intervalles ouverts disjoints deux à deux.
\end{exo}

\subsection{Image continue d'un connexe par arcs}

On dispose d'un version plus générale du théorème des valeurs intermédiaires.

\begin{prop}%prop
	L'image d'un connexe par arcs par une application continue est un connexe par arcs.
\end{prop}

\begin{exo}
	Soit $E$ un $\K$-espace vectoriel normé (éventuellement de dimension infinie). La sphère unité de $E$ est connexe par arcs sauf si $\K=\R$ et $\dim(E)=1$.
\end{exo}

\begin{attention}
	L'image réciproque d'un connexe par arcs par une application continue peut très bien ne pas être connexe par arcs. Produire un exemple.
\end{attention}

\begin{cor}%cor
	Si $f:E\to\R$ est continue, l'image de tout connexe par arcs de $E$ est un intervalle.
\end{cor}

\begin{ex}
	\begin{enumerate}
		\item $\R$ et $\C$ ne sont pas homéomorphes.
		\item Une injection continue d'un intervalle de $\R$ dans $\R$ est strictement monotone.
	\end{enumerate}
\end{ex}

\begin{exo}
	Montrer qu'il n'existe pas de bijection continue de $[0,1]^2$ sur $[0,1]$.
\end{exo}

\begin{rmq}
	En utilisant le deuxième exemple qui précède, on en déduit qu'il n'existe pas non plus de bijection continue de $[0,1]$ sur $[0,1]^2$.
\end{rmq}

\begin{exo}
	\begin{enumerate}
		\item Montrer que $\mathcal{GL}_n(\R)$ n'est pas connexe par arcs. %déterminant
		\item Quelles sont ses composantes connexes par arcs ? %$\mathcal{GL}_n^-(\R)$ et $\mathcal{GL}_n^+(\R)$
		\item Qu'en est-il de $\mathcal{GL}_n(\C)$ ?
	\end{enumerate}
\end{exo}

\subsection{Connexité (HP)}

\begin{dfn}
	Un espace topologique $X$ est \textbf{connexe} si les seules parties de $X$ à la fois ouvertes et fermées sont $X$ et $\emptyset$.
\end{dfn}

\begin{exo}
	Dans un connexe, quelles sont les parties de frontière vide ?
\end{exo}

\begin{prop}%prop
	Tout connexe par arcs est connexe.
\end{prop}
\begin{proof}
	La fonction indicatrice d'une partie $P$ ouverte et fermée est continue.
\end{proof}

\begin{dfn}
	On dit qu'une application $f:A\to F$ est localement constante sur $A$ si elle constante au voisinage de tout point.
\end{dfn}

\begin{cor}%cor
	Une application localement constante sur un connexe (et donc sur un connexe par arcs) est constante.
\end{cor}

\begin{attention}
	Il existe des connexes non connexes par arcs : voir l'exercice Moulin.Comp.39.
\end{attention}

\begin{exo}
	En utilisant les résultats de l'exercice sur les composantes connexes par arcs d'un ouvert d'un espace vectoriel normé, montrer qu'un ouvert connexe d'un espace vectoriel normé est connexe par arcs.
\end{exo}

\section{Espaces vectoriels normés de dimension finie}

Soit $E$ un $\K$-espace vectoriel de dimension finie non nulle (le cas $E=\{0\}$ est trivial).

\subsection{Équivalence des normes}

Si $\mathcal{B}=(e_1,\ldots,e_n)$ est une base de $E$, on peut définie la norme $N_\infty$ sur $E$ par : $$N_\infty\left(\displaystyle\sum_{i=1}^{n}x_ie_i\right)=\underset{i\in\llbracket1,n\rrbracket}{\max}\lvert x_i\rvert.$$
L'isomorphisme $f:(x_i)\mapsto\displaystyle\sum_{i=1}^{n}x_ie_i$ est une \textit{isométrie} de $(\K^n,\lVert\cdot\rVert_\infty)$ sur $(E,N_\infty)$, c'est-à-dire que l'on a : 
$$\forall x\in\K^n,\ N_\infty(f(x))=\lVert x\rVert_\infty.$$

Elle est donc continue, ainsi que sa réciproque. On en déduit :

\begin{lem}%lem
	Les compacts de $(E,N_\infty)$ sont les fermés bornés. \\
	En particulier, la sphère unité de $(E,N_\infty)$ est compacte.
\end{lem}

\begin{thm}%thm
	Sur un espace $E$ de dimension finie, toutes les normes sont équivalentes.
\end{thm}

\begin{proof}
	Soit $N$ une norme sur $E$.
	\begin{itemize}
		\item On montre que $N$ est dominée par $N_\infty$, ce qui prouve que $N$ est continue sur $(E,N_\infty)$.
		\item Donc $N$ est minorée sur la sphère unité de $(E,N_\infty)$ par un réel strictement positif.
	\end{itemize}
\end{proof}

\begin{rmq}
	Il en résulte qu'en dimension fini, pour les notions de limite, de continuité, de continuité uniforme, d'application lipschitzienne, le choix de la norme est indifférent. Le plus souvent, il n'est alors pas précisé.
\end{rmq}

\subsection{Compacité en dimension finie}

\begin{thm}%thm
	Les parties compactes d'un espace vectoriel normé de dimension finie sont exactement ses parties fermées bornées.
\end{thm}

Exercice Moulin.Comp.49

\begin{cor}%cor
	\begin{enumerate}
		\item Toute suite bornée d'un espace vectoriel de dimension finie admet une valeur d'adhérence.
		\item Toute suite bornée d'un espace vectoriel de dimension finie admettant une unique valeur d'adhérence est convergente.
	\end{enumerate}
\end{cor}

\begin{prop}%prop
	Soit $F$ un sous-espace vectoriel de dimension finie d'un espace vectoriel normé $E$. Alors $F$ est fermé dans $E$.
\end{prop}

\begin{rmq}
	Ce qui est remarquable dans ce résultat, c'est que $F$ est fermé pour n'importe quelle norme de $E$. D'ailleurs, toutes les normes de $E$ induisent la même topologie sur $F$.
\end{rmq}

\begin{prop}[HP]%prop
	Dans chacun des deux cas suivants :
	\begin{itemize}
		\item $F$ est un sous-espace vectoriel de dimension finie d'un espace vectoriel normé $E$ ;
		\item $F$ est une partie fermée non vide d'un espace vectoriel $E$ de dimension finie ;
	\end{itemize}
	la distance à $F$ de tout point $x\in E$ est atteinte en au moins un point de $F$
\end{prop}
\begin{proof}
	On considère le compact $K=F\cap B_f(x,d(a,x))$ où $a\in F$.
\end{proof}

\begin{exo}
	Donner un exemple de partie fermée $F$ d'un espace vectoriel normé et d'un élément $a\in E$ dont la distante à $F$ est non atteinte.
\end{exo}

\begin{exo}
	Soit $F$ un sous-espace vectoriel strict de dimension finie d'un espace vectoriel normé $E$. Montrer qu'il existe une élément $x\in E$ tel que $\lVert x\rVert=1$ et $d(x,F)=1$.
\end{exo}

\begin{rmq}
	La boule unité fermée d'un espace vectoriel de dimension finie est compacte puisque fermée bornée en dimension finie. Nous avons vu au début du chapitre que la boule unité fermée pouvait être non compacte. Pire :
\end{rmq}

\begin{thm}[Théorème de Riesz (HP)]%thm
	Soit $E$ un espace vectoriel normé dont la boule unité $B$ est compacte. Alors $E$ est de dimension finie.
\end{thm}
\begin{proof}
	On suppose que $E$ n'est pas de dimension finie, et l'on construit dans $E$, à l'aide de l'exercice précédent, une suite $(x_n,F_n)$ avec : $$F_n=\Vect(x_1,\ldots,x_n)\ \ \text{et}\ \ x_{n+1}\in B\ \ \text{tel que}\ \ d(x_{n+1},F_n)=1.$$ La suite $(x_n)$ est à valeurs dans $B$ et ne possède pas de valeur d'adhérence.
\end{proof}

\subsection{Continuité des applications (multi)linéaires et polynomiales}

\begin{thm}%thm
	Soit $E$ un espace vectoriel normé de dimension finie et $F$ un espace vectoriel normée de dimension quelconque. Toute application linéaire de $E$ dans $F$ est continue.
\end{thm}

\begin{rmq}
	Les applications affines sur un espace vectoriel de dimension finie sont donc continues comme somme d'une constante et d'une application linéaire.
\end{rmq}

\begin{ex}
	Si $u\in\mathcal{L}(E,F)$, où $E$ est de dimension finie, alors il existe $x\in E$ non nul tel que $\lVert u(x)\rVert=\lVert u\rVert_{op}\lVert x\rVert$. 
\end{ex}

En particulier, les formes linéaires coordonnées dans la base $\mathcal{B}$ sont continues. Par combinaison linéaire et produit :

\begin{prop}%prop
	Toute application d'un $\K$-espace vectoriel de dimension finie polynomiale en les coordonnées dans la base $\mathcal{B}$ est continue sur $E$.
\end{prop}

\begin{ex}
	L'application déterminant $\det:\mathcal{M}_n(\K)\to\K$ est une application continue.
\end{ex}

\begin{prop}%prop
	Soit $E_1,\ldots,E_p$ des $\K$-espaces vectoriels de dimension finie. Toute application $p$ linéaire de $E_1\times\cdots\times E_p$ dans un espace vectoriel $F$ est continue.
\end{prop}
\begin{proof}
	\begin{itemize}
		\item Si $F=\K$, une telle application est polynomiale.
		\item Si $F$ est de dimension finie, on utilise les fonctions coordonnées.
		\item Dans le cas général, on se ramène à un sous-espace vectoriel de $F$ de dimension finie.
	\end{itemize}
\end{proof}

\begin{ex}
	Si $E$ est un espace vectoriel de dimension finie $n$ muni d'une base $\mathcal{B}$, l'application $(u_1,\ldots,u_n)\mapsto\det_\mathcal{B}(u_1,\ldots,u_n)$ est continue sur $E^n$.
\end{ex}

Exercice Moulin.Comp.53 (première question).

\begin{cor}%cor
	Soit $E_1,\ldots,E_p$ des $\K$-espaces vectoriels normés de dimension finie et $f$ une application $p$-linéaire de $e_1\times\cdots\times E_p$ dans un espace vectoriel $F$. Alors il existe un réel $C$ tel que : $$\forall (_1,\ldots,x_p)\in E_1\times\cdots\times E_p,\ \lVert f(x_1,\ldots,x_p)\rVert\leqslant C\lVert x_1\rVert\cdots\lVert x_p\rVert.$$
\end{cor}

\begin{exo}
	Soit $E$ une $\K$-algèbre de dimension finie. Montrer qu'il existe une norme sur $E$ qui soit sous-multiplicative, c'est-à-dire vérifiant : $$\forall (x,y)\in E^2,\ \lVert xy\rVert\leqslant\lVert x\rVert\lVert y\rVert.$$
\end{exo}

\subsection{Fonctions coordonnées}

Soit $f:A\subset E\to F$. Si $f$ est de dimension finie muni d'une base $\mathcal{B}=(e_1,\ldots,e_n)$, on peut écrire : $$f(x)=\sum_{i=1}^{n}f_i(x)e_i$$ où les $f_i$ sont les fonctions coordonnées dans la base $\mathcal{B}$. Soit $k\in\llbracket1,n\rrbracket$. On note $p_k$ la $k$ème forme coordonnée dans la base $\mathcal{B}$ et $\begin{array}{lrcl}
	q_k:&\K&\longrightarrow& F\\
	&t&\longmapsto& te_k.
\end{array}$.
On a alors $f_k=p_k\circ f$ et $f=\displaystyle\sum q_i\circ f_i$.

\begin{met}
	\begin{itemize}
		\item $f$ admet une limite finie en $a$ adhérent à $A$ si, et seulement si, $f_i$ admet une limite finie pour tout $i\in\llbracket1,n\rrbracket$.
		\item $f$ est continue sur $A$ si, et seulement si, toutes les fonctions $f_i$ sont continues sur $A$.
		\item On a le même résultat en remplaçant "continue" par "uniformément continue" ou "lipschitzienne".
	\end{itemize}
\end{met}

\begin{rmq}
	En particulier, une suite $(u_k)_{k\in\N}=\displaystyle\left(\sum_{i=1}^{n}u_{i,k}e_i\right)$ à valeurs dans $F$ converge si, et seulement si, toutes les suites coordonnées $(u_{i,k})_{k\in\N}$ convergent.
\end{rmq}

\begin{thm}%thm
	Dans un espace vectoriel normé de dimension finie, toute série absolument convergente est convergente.
\end{thm}
\begin{proof}
	Le faire composante par composante.
\end{proof}

\end{document}